\documentclass{mycv}

\name[ZHANG Xinyi]{ZHANG Xinyi}
\address{DLB 625G, 34 Renfrew Rd, HKBU, Hong Kong}
\email{csxyzhang@comp.hkbu.edu.hk}
\phone{+852-56027517}
\homepage{https://linekm.github.io}
% \github{xu-cheng}
\linkedin{xinyi-zhang-hkbu}

\begin{document}

\maketitle%

\begin{summary}
Postdoctoral Fellow in the Department of Computer Science at Hong Kong Baptist University with experience in designing and implementing various high-performance, secure database and distribution-aware data structures.
\end{summary}

\section{Education}

\subsection{Hong Kong Baptist University}[Hong Kong, China]
\vspace{-\parskip}%
\begin{itemize}[label={}]
  \item Ph.D.\ in Computer Science \printdate{Sep 2021~--~Jul 2026}
  \item Supervisor: \href{https://www.comp.hkbu.edu.hk/~xujl}{Prof.~Xu Jianliang}
\end{itemize}

\subsection{Huazhong University of Science and Technology}[Wuhan, China]
\vspace{-\parskip}%
\begin{itemize}[label={}]
  \item Bachelor of Engineering in Electronics \& Information Engineering \printdate{Sep 2017~--~Jun 2021}
  \item \textbf{GPA}: 3.7/4.0 (Rank: 15\%)
\end{itemize}

\section{Professional \\ Experience}

\subsection{Hong Kong Baptist University}[Hong Kong, China]
\begin{positions}
  \entry{Postdoctoral Fellow}{Aug 2026~--~Present}
  \entry{Ph.D.\ Candidate}{Sep 2021~--~Jul 2026}
\end{positions}

\begin{itemize}
  \item Supervisor: \href{https://www.comp.hkbu.edu.hk/~xujl}{Prof.~Xu Jianliang}
  \item Designed novel algorithms and indexes for relational and vector databases to support efficient and secure query processing in a wide range of systems.
  \item Designed novel distribution-aware data structures to handle complex database workloads while ensuring robust performance.
  \item Research resulted in publications at top-tier conferences.
\end{itemize}

\subsection{Nanyang Technological University}[Singapore]
\begin{positions}
  \entry{Visiting Scholar}{March 2026~--~June 2026}
\end{positions}
\begin{itemize}
  \item Supervisor: \href{https://qichen-wang.github.io}{Prof.~Wang Qichen}
  \item Focus on "Learned Indexes for Query Optimization," specifically aiming to integrate learned index structures into structure-guided execution plans for acyclic queries.
\end{itemize}

\subsection{Huazhong University of Science and Technology}[Wuhan, China]
\begin{positions}
  \entry{Researcher}{May 2019~--~Feb 2020}
\end{positions}

\begin{itemize}
  \item Design a fast consensus algorithm of blockchain combining SOTA consensus algorithm.
  \item Build a blockchain system for ASTRI as a integrity-assured record for distributed database.
\end{itemize}

% \subsection{Huazhong University of Science and Technology}[Wuhan, China]
% \begin{positions}
%   \entry{Core Member}{Oct 2019~--~Apr 2021}
% \end{positions}

% \begin{itemize}
%   \item Research and implement the Algorand consensus algorithm to verify its feasibility for ASTRI's BFC 2.0 system.
% \end{itemize}

\section{Research \\ Interests}

\begin{itemize}
  \item Oblivious query processing for data federation and outsourcing cloud computing.
  \item Scalable secure query processing for vector database based on Trusted Execution Environments.
  \item Robust updatable learned index under real-world workload.
  \item Distributed vector database design
\end{itemize}

\section{Publications}%

% \textbf{Complete List:}
% \href{https://scholar.google.com/citations?user=DKG_JaAAAAAJ}{Google Scholar \textsf{\footnotesize [DKG\_JaAAAAAJ]}}%
% {~~$\cdot$~~}%
% \href{https://dblp.org/pers/hd/x/Xu_0004:Cheng}{DBLP \textsf{\footnotesize [Xu\_0004:Cheng]}}%

\publications[keyword={selected}]{publications.bib}

{
\footnotesize%
\textsuperscript{\textdagger}These authors contributed equally.
}

\section{Research \\ Projects}
\subsection{FedKNN: Secure Federated k-Nearest Neighbor Search}[\textbf{SIGMOD '24}]
\begin{itemize}
  \item Proposed FedKNN, a system for secure and privacy-preserving federated kNN search, achieving up to 4.8$\times$ efficiency improvement over SOTA by optimizing local computations with a Distribution-Aware algorithm.
\end{itemize}

\vspace{+5pt}

\subsection{HIRE: A Hybrid Learned Index for Robust and Efficient Performance}[\textbf{SIGMOD '26}]
\begin{itemize}
  \item Designed HIRE, a hybrid in-memory index combining learned predictions with traditional robustness, which reduces tail latency by 98\% and improves throughput by 41.7× compared to SOTA approaches.
\end{itemize}

\subsection{Privacy-Preserving Approximate Nearest Neighbor Search}[\textbf{Under Review}]
\begin{itemize}
  \item Proposed a secure ANN index structure within trusted execution environments, mitigating side-channel attacks by enforcing oblivious memory access patterns on proximity graph indexes. In this paper, we achieve plaintext-level throughput with provable privacy guarantees by devising a distribution-aware oblivious access algorithm, significantly outperforming SOTA solutions.
\end{itemize}

\subsection{Authenticated Approximate Nearest Neighbor Search}[\textbf{Under Revision}]
\begin{itemize}
  \item Designed VIDA, a novel scheme for authenticated graph-based ANN search. VIDA leverages vector commitments and Inner Product Arguments to verify distance computations, thereby drastically reducing the overall data transmission size. To further optimize performance, we introduce MeST, a proof aggregation mechanism that compresses the verification of multiple vectors into a single cryptographic proof. In extensive experiment, our solution achieves up to a 5.2x speedup in end-to-end latency and reduces data transfer time by 99.3\%.
\end{itemize}

\section{Talks}

\begin{enumerate}
  \item FedKNN: Secure Federated k-Nearest Neighbor Search, \emph{2024 ACM SIGMOD International Conference on Management of Data}, Santiago, Chile, Jun.~2024.
  
  \item Secure Federated kNN Search, \emph{Huawei}, Online, Nov.~2024.

  \item HIRE: A Hybrid Learned Index for Robust and Efficient Performance under Mixed Workloads, \emph{2026 ACM SIGMOD International Conference on Management of Data}, Bengaluru, India, Jun.~2026.
\end{enumerate}

\section{Services}
\begin{itemize}
    \item Reviewer, IEEE Transactions on Parallel and Distributed Systems (TPDS) \printdate{2025}
    \item Reviewer, IEEE Transactions on Knowledge and Data Engineering (TKDE) \printdate{2024-2026}
    \item Reviewer, Springer Data Science and Engineering (DSE) \printdate{2024}
    \item External Reviewer, ACM International Conference on Management of Data (SIGMOD) \printdate{2025}
    \item External Reviewer, International Conference on Very Large Data Bases (VLDB) \printdate{2023--2025}
    \item External Reviewer, IEEE International Conference on Data Engineering (ICDE) \printdate{2022--2025}
\end{itemize}

\section{Skills}

\begin{description}
  \item[Programming] Rust, C/C++, Golang, Python, \LaTeX, Bash
  \item[Tools] Docker, Vim, Git, Linux
  \item[Languages] English, Chinese
\end{description}

\section{Awards}

\begin{itemize}
  \item Academic Excellence Scholarship, Huazhong University of Science and Technology \printdate{2018,~2021}
  % \item First Prize of ADI SDR Cup, Huazhong University of Science and Technology \& ADI \printdate{2019}
  \item Excellent Teaching Assistant Performance Award, Hong Kong Baptist University \printdate{2022,~2023}
  \item RPg Research Performance Award, Hong Kong Baptist University \printdate{2024,~2025}
  \item ACM SIGMOD Student Travel Award, ACM \printdate{2024,~2026}
  \item VLDB Student Travel Award, VLDB Endowment \printdate{2024}
\end{itemize}

\section{Referees}
\begin{itemize}
  \item \href{https://www.comp.hkbu.edu.hk/~xujl}{Prof.~Xu Jianliang}, Head \& Chair Professor, Hong Kong Baptist University
  \item \href{https://qichen-wang.github.io}{Dr.~Wang Qichen}, Assistant Professor, Nanyang Technological University
  \item \href{https://www.comp.hkbu.edu.hk/~bchoi/}{Prof.~Byron Choi}, Associate Head \& Professor, Hong Kong Baptist University
\end{itemize}

\end{document}
